\documentclass[12pt]{article}

% --- Encoding / language (XeLaTeX-safe) ---
\usepackage{fontspec}
\usepackage[english, russian]{babel}
\defaultfontfeatures{Ligatures=TeX}
\setmainfont{Times New Roman}

% --- Math / theorem ---
\usepackage{amsmath}
\usepackage{amssymb}
\usepackage{amsthm}
\usepackage{stmaryrd}
\usepackage{breqn}
\usepackage{euscript}
\usepackage{mathrsfs}
\usepackage{enumerate}
\usepackage{empheq}

% --- Graphics / layout ---
\usepackage{graphicx}
\usepackage{wrapfig}
\usepackage{subcaption}
\usepackage{float}
\usepackage{xcolor}
\usepackage{geometry}

\geometry{
    left=2.5cm,
    right=2.5cm,
    top=2cm,
    bottom=2cm,
    bindingoffset=0cm
}

% --- Extra utilities ---
\usepackage{pifont}
\usepackage{lipsum}
\usepackage{todonotes}
\usepackage{csquotes}

% --- Bibliography ---
\usepackage[backend=biber,style=alphabetic]{biblatex}
\addbibresource{references.bib}

% --- Hyperref (must be loaded last) ---
\usepackage{hyperref}

% --- Custom imports ---
\input{notations}

% --- Safe image inclusion (shows empty figure if file missing) ---
\newcommand{\safeincludegraphics}[2][]{%
    \IfFileExists{#2}{%
        \includegraphics[#1]{#2}%
    }{%
        \textcolor{gray}{\rule{0.8\linewidth}{0.6\linewidth}}%
    }%
}

% --- Theorems ---
\newtheorem{proposition}{Proposition}
\newtheorem{theorem}{Theorem}
\theoremstyle{remark}
\newtheorem{remark}{Remark}

% --- Colors ---
\definecolor{darkred}{RGB}{139,0,0}

% --- Date helper ---
\newcommand{\todayYYMMDD}{\the\year-\the\month-\the\day}


\begin{document}

\hfill
% \todayYYMMDD
\begin{center}
% \section*{Custom course or document title}
\textit{eva hw2 задание1}\\
% \textit{Optional subtitle}\\
\textit{Author: Полозков Дмитрий}\\
% \textit{Supervisor: Name}\\
\textit{Date: \today}
\end{center}

% \tableofcontents
% \newpage 

\section*{Задание 1}

\begin{enumerate}
    \item 
    В первую очередь, заметим симметричность копулы по $u_1$, $u_2$.
    Таким образом, любое свойство, верное для $u_1$ автоматически верно для $u_2$.

    Первое св-во копулы: равенство 0 когда один из аргументов равен 0.
    У нас: $C(0,\,u_2\;\;\theta)=0+0+0$ т.к. первое слагаемое имеет множитель $\min(0,\,u_2)=0$; второй - домножение на $u_1=0$; трерий - множитель $\max(0+u_2-1,\,0)=0$.

    Второе св-во: $C$ не убывает по $u_1$. 
    Также выполнено т.к. $\min$ от неуб. аргументов - неуб.; $u_1u_2$ - неуб. функция; $\max$ от неуб. арг. - неуб.

    Третье св-во: 
    \begin{align*}
        C(1,\,u_2;\;\theta)&=\frac{\theta^2(1+\theta)}{2}u_2+(1-\theta^2)u_2+\frac{\theta^2(1-\theta)}{2}u_2\\
        &=\frac{2\theta^2}{2}u_2+(1-\theta^2)u_2\\
        &=u_2.
    \end{align*}
    \item 
    По т. Склара, 
    \[
        \P(X_1\le x_1,\,X_2\le x_2)=C(F_1(x_1),\,F_2(x_2);\;\theta).
    \]
    Значит, по формуле включений исключений,
    \begin{align*}
        \P(X_1>10,\,X_2>2)&=1-\P(X_1<10)-\P(X_2<2)+\P(X_1<10,\,X_2<2)\\
        &=1-F_1(10)-F_2(2)+C(F_1(10),\,F_2(2);\;\theta). 
    \end{align*}
    Осталось только подставить значения и получить функцию от $\theta$
    
    :)
    \item
    Заметим, независимость эквивалентна выполнению
    \[
        \P(X_1\le x_1,\,X_2\le x_2)=\P(X_1\le x_1)\P(X_2\le x_2)
    \]
    для любых $x_1,\,x_2$.
    Данное можно переписать как
    \begin{equation}
        \label{eq:indep}
        C(F_1(x_1),\,F_2(x_2);\;\theta)=F_1(x_1)F_2(x_2).
    \end{equation}
    Заметим, при $\theta=0$ имеем 
    \[
        C(u_1,\,u_2;\;\theta)=0+(1-0)u_1u_2+0=u_1u_2,
    \]
    что дает \eqref{eq:indep}.
\end{enumerate}

\end{document}

